\section{Periods and Riemann Bilinear Relations}
\label{sec:periods}

\begin{layman}
A \emph{period} is a number you get by integrating a holomorphic
1-form along a closed loop on a Riemann surface. For an elliptic
curve (genus 1, topologically a doughnut) there are two
fundamental loops — once through the hole, once around it — and
correspondingly two fundamental periods. They are complex numbers,
typically transcendental, and historically they're where the entire
theory began: the \emph{elliptic integrals} of the eighteenth
century are exactly the period functions of genus-one curves.
Riemann's insight was to abstract the period across all genera:
the complete \emph{period pairing} of a Riemann surface is the
linear map sending each homology class to ``the linear functional
\(\omega \mapsto \int \omega\) along that class.''

The first major theorem of this section says periods only depend on
the homology class of the loop, not the loop itself. The mechanism
is Stokes' theorem: if two loops are homologous they bound a
surface region, the integrand is closed (because holomorphic
1-forms are closed), and Stokes turns the region's surface integral
into the difference of the loop integrals — which is zero. This
homology invariance is what lets us call ``\(\Pi_X\)'' a
\emph{pairing} on \(H_1(X,\Z)\) at all.

The deeper theorem of this section is the \emph{Riemann bilinear
relations}. They identify a wedge integral of two 1-forms over the
whole surface with a sum of period products picked out by a
\emph{symplectic basis} of homology — a normalised basis where the
\(g\) loops paired with \(g\) other loops cross each other exactly
once and don't otherwise interact. The proof is a beautiful
choreography: cut the surface along the symplectic basis to lay it
flat as the \(4g\)-gon of the previous section; on the polygon's
flat interior, find a primitive \(F\) of one form (using simple
connectedness); apply Stokes on the polygon, where the boundary
collapses by the gluing pattern into a sum over symplectic-basis
periods. From this identity plus a positivity statement
(\(i\int \omega \wedge \bar\omega > 0\) for nonzero \(\omega\)) we
conclude that the \(2g\) symplectic-basis period vectors are
linearly independent over the reals.

Why bother? Linear independence over the reals means the integer
combinations of period vectors form a \emph{full lattice} — a
discrete grid spanning the entire real vector space underlying the
dual of the 1-forms. That full lattice is the \emph{period lattice};
the next section quotients by it to build the Jacobian. Every
geometric and analytic property of the Jacobian — compactness,
manifold structure, group structure, complex structure — flows
from the period lattice being full. Without these bilinear
relations there would be no Jacobian.
\end{layman}

For a singular one-cycle \(\gamma\) and \(\omega\in\HX\), define
\[
  \Pi_X(\gamma)(\omega)=\int_{\gamma}\omega.
\]
Since holomorphic one-forms are closed, the integral only depends on the
homology class of \(\gamma\). This gives the period pairing of
Definition~\ref{def:period-pairing}.

\subsection{Homology invariance of periods}

\begin{lemma}[Homology invariance of periods]
\label{lem:period-homology-invariance}
\uses{thm:stokes-on-rs-with-boundary,lem:holomorphic-form-is-closed}
\notready
If \(\gamma\) and \(\gamma'\) represent the same class in \(H_1(X,\Z)\), then
\[
  \int_{\gamma}\omega=\int_{\gamma'}\omega
\]
for every \(\omega\in\HX\).
\lean{JacobianChallenge.Periods.period_homology_invariance_statement}
\end{lemma}

\begin{layman}
Pick a holomorphic 1-form on the surface — a tiny linear gadget at every
point — and a closed loop. Integrating the form along the loop produces a
single complex number, the \emph{period}. The lemma says: if you wiggle
the loop continuously, dragging it through the surface, the period
doesn't change. More precisely, two loops that represent the same class
in homology — meaning they bound a 2-dimensional region between them —
give the same period. The argument is Stokes' theorem applied to the
region enclosed, plus the fact that holomorphic 1-forms have zero
exterior derivative. Why bother? This makes the period a function of the
\emph{homology class} of a loop, not the loop itself. Since the homology
of a compact genus-\(g\) surface is a small free abelian group of rank
\(2g\), we now have a clean finite-dimensional ``period pairing''
between homology and 1-forms — the ingredient that builds the Jacobian's
lattice.
\end{layman}

\begin{proof}
The difference \(\gamma-\gamma'\) is a boundary. Holomorphic one-forms are
closed (Lemma~\ref{lem:holomorphic-form-is-closed}). Stokes' theorem
(Theorem~\ref{thm:stokes-on-rs-with-boundary}) gives that the integral of
a closed form over a boundary vanishes.
\end{proof}

\begin{lemma}[Holomorphic 1-forms are closed]
\label{lem:holomorphic-form-is-closed}
For \(\omega\in\HX\), \(d\omega=0\).
\depends{Local computation: \(d(h(z)\,dz)=\frac{\partial h}{\partial \bar
z}\,d\bar z\wedge dz=0\) by Cauchy-Riemann; globalization via the chart
atlas.}
\lean{JacobianChallenge.Blueprint.holomorphic_form_is_closed}
\leanok
\end{lemma}
\begin{proof}\leanok\end{proof}

\begin{layman}
A holomorphic 1-form, written locally as a holomorphic function times the
differential of the coordinate, has zero exterior derivative. The
exterior derivative of \(h(z)\,dz\) — the gadget that measures
``infinitesimal swirl'' — vanishes because of the Cauchy–Riemann
equations: holomorphic functions don't depend on the conjugate
coordinate, so the only term that could appear cancels automatically.
Concretely, holomorphic 1-forms are \emph{closed}: no swirl, no
infinitesimal failure of being a total derivative. Why bother? Closed
1-forms are exactly the gadgets where Stokes' theorem makes the integral
over a boundary equal to zero. That fact is what makes the period
computation depend only on the homology class of a loop, and what
underlies all the rigidity statements about the period lattice further
down.
\end{layman}

\begin{theorem}[Stokes' theorem on Riemann surfaces with boundary]
\label{thm:stokes-on-rs-with-boundary}
For an oriented compact Riemann surface with boundary \(M\) and a smooth
1-form \(\omega\) on \(M\),
\[
  \int_{\partial M}\omega = \int_{M} d\omega.
\]
\depends{ABSENT in Mathlib v4.28.0 (Mathlib has Stokes only on box /
rectangle subsets of \(\R^n\), not on manifolds with boundary). This is
the single largest classical-analysis gap for the period-lattice route
per the project inventory.}
\lean{JacobianChallenge.Blueprint.Sec03.stokes_on_rs_with_boundary}
\leanok
\end{theorem}

\begin{layman}
Stokes' theorem is the grandfather of all integration-by-parts results.
On a compact 2-dimensional surface with boundary — think of a disk, or
a piece cut out of a donut — and any smooth 1-form on it, the integral
of the form along the boundary equals the integral of its exterior
derivative over the surface. ``Boundary integral'' walks you along an
edge; ``derivative-of-form integrated over the surface'' is a 2D
integral. The theorem says these two flavors of integration are
miraculously equal. Mathlib only has this for boxes in Euclidean space,
not for general curved manifolds with boundary, which makes it the
single largest classical-analysis gap on our path. Why bother? Stokes
is the engine that turns local differential properties into global
integral identities. The entire Riemann bilinear story — and homology
invariance of periods — cashes out via Stokes on the polygonal model.
\end{layman}

\subsection{Riemann bilinear relations and full real rank}

The Riemann bilinear relations come from cutting \(X\) along a
symplectic basis to obtain the standard \(4g\)-gon
(Theorem~\ref{thm:polygonal-model}) and applying Stokes there with
the help of the polygon-interior primitive
(Lemma~\ref{lem:primitive-on-polygon}).

\begin{definition}[Symplectic basis of \(H_1(X,\Z)\)]
\label{def:symplectic-basis}
A symplectic basis of \(H_1(X,\Z)\) is a \(\Z\)-basis
\(a_1,\dots,a_g,b_1,\dots,b_g\) for which the algebraic intersection
pairing satisfies
\((a_i,a_j)=(b_i,b_j)=0\) and \((a_i,b_j)=\delta_{ij}\).
\depends{Mathlib has the algebraic theory of symplectic bilinear forms;
the existence of a symplectic basis on a compact orientable surface is
a topological theorem (every nondegenerate alternating form on a free
\(\Z\)-module of rank \(2g\) admits a symplectic basis). The
polygonal model of Section~\ref{sec:polygonal-model} provides a
direct construction: the boundary identification of
\(\mathrm{Polygon4g}(g)\) reads off the symplectic generators
\(a_i, b_i\) by handle index.}
\lean{JacobianChallenge.Blueprint.Sec03.isSymplecticBasis}
\leanok
\end{definition}

\begin{layman}
On a compact genus-\(g\) surface, the closed-loop equivalence classes
form a free abelian group of rank \(2g\): exactly \(2g\) independent
``ways to go around.'' On a donut (\(g = 1\)) there are two: once
around the hole, and once through it. On a two-handled surface there
are four: two per handle. A \emph{symplectic basis} chooses \(g\) of
those loops to be ``\(a\)'s'' and the other \(g\) to be ``\(b\)'s,''
arranged so that any \(a\) loop crosses its matching \(b\) loop exactly
once and crosses every other basis loop zero times. It's a sharply
normalized basis where the geometric intersection matrix is the
simplest possible: zeros and ones in a tidy pattern. Why bother? The
Riemann bilinear relations and the period-lattice argument become
genuinely tractable only after this normalization — a tangle of loops
becomes a tidy \(2g\)-vector with predictable interactions.
\end{layman}

\begin{theorem}[Bilinear identity from Stokes on the polygon]
\label{thm:bilinear-from-stokes}
\uses{thm:polygonal-model,lem:primitive-on-polygon,thm:stokes-on-rs-with-boundary}
For holomorphic one-forms \(\omega,\eta\) on \(X\),
\[
  \int_X \omega\wedge\eta
  =
  \sum_{i=1}^{g}
  \left(
    \int_{a_i}\omega\int_{b_i}\eta
    -
    \int_{b_i}\omega\int_{a_i}\eta
  \right).
\]
\lean{JacobianChallenge.Blueprint.bilinear_from_stokes}
\notready
\end{theorem}

\begin{layman}
Take two holomorphic 1-forms on the surface. The integral of their
wedge product (a 2-form) over the whole surface miraculously equals a
finite sum of products of single-loop periods: pair up the symplectic
basis cycles \((a_i, b_i)\), compute four numbers (the integrals of
each form along each cycle), multiply across the pairs in a specific
alternating way, and add. The bookkeeping mirrors the simple algebraic
``determinant'' of a \(2 \times 2\) period matrix, summed over handles.
The proof is Stokes on the polygonal model: turn the surface integral
into a polygon-edge integral via the primitive of one of the forms,
then use the gluing pattern to fold opposite edges together — the
difference of integrand values across a glued pair is exactly a
period. Why bother? This identity is the foundation of the Riemann
bilinear relations and (combined with positivity in the next step)
forces the period lattice to be \emph{full}.
\end{layman}

\begin{proofsketch}
Pull \(\omega,\eta\) back to the polygon \(P\). On \(P\), let \(F\) be a
primitive of \(\omega\) (Lemma~\ref{lem:primitive-on-polygon}). Then
\[
  \omega\wedge\eta = d(F\eta).
\]
Apply Stokes to integrate over \(\partial P\). The side identifications
group the boundary integral into pairs of cycles; the algebraic
combination of \(F\) values across paired sides is exactly the symplectic
bilinear sum on the right.
\end{proofsketch}

\begin{theorem}[Hermitian positivity \(i\int_X\omega\wedge\bar\omega>0\)]
\label{thm:hermitian-positivity}
\uses{thm:bilinear-from-stokes}
For a nonzero holomorphic 1-form \(\omega\),
\[
  i\int_X \omega\wedge\overline{\omega} > 0.
\]
\lean{JacobianChallenge.Blueprint.hermitian_positivity}
\notready
\end{theorem}

\begin{layman}
For any nonzero holomorphic 1-form, the integral of the form wedged with
its complex conjugate (multiplied by an orientation-fixing factor of
\(i\)) is strictly positive. Locally, in a chart where the form looks
like \(h(z)\,dz\), the integrand simplifies to \(2|h(z)|^2\) times the
ordinary area element — a nonnegative density that's strictly positive
wherever \(h\) doesn't vanish. Integrate this nonnegative thing over
the whole surface: the result is positive unless the form vanished
identically, which it didn't by assumption. The stray \(i\) at the
front handles the orientation factor; without it you'd get a strictly
negative number instead. Why bother? This positivity, paired with the
bilinear identity from the previous step, is what forces the period
vectors to be linearly independent over the real numbers. No
positivity, no lattice.
\end{layman}

\begin{proofsketch}
Locally \(\omega = h(z)\,dz\), so
\(i\,\omega\wedge\bar\omega = 2|h(z)|^{2}\,dx\wedge dy\),
which is a nonnegative volume form, strictly positive at any point
where \(\omega\) is nonzero. Integrate over \(X\) and use connectedness
to conclude positivity.
\end{proofsketch}

\begin{theorem}[Period vectors have full real rank]
\label{thm:period-vectors-full-real-rank}
\uses{thm:bilinear-from-stokes,thm:hermitian-positivity,def:symplectic-basis}
The \(2g\) period vectors
\(\Pi_X(a_1),\dots,\Pi_X(a_g),\Pi_X(b_1),\dots,\Pi_X(b_g)
   \in\HdualX\)
are linearly independent over \(\R\).
\lean{JacobianChallenge.Periods.period_vectors_linearIndependent_of_symplectic}
\leanok
\end{theorem}

\begin{layman}
A compact genus-\(g\) surface has \(2g\) independent loops in homology
(the symplectic basis). Each loop, integrated against holomorphic
1-forms, gives a vector of length \(g\) in the dual of the 1-form
space — the \emph{period} of that loop. Bundle the \(2g\) such
vectors. The theorem: they are linearly independent over the real
numbers. So inside a \(2g\)-real-dimensional space we have \(2g\)
independent vectors; they form a basis. The proof is a beautiful
mash-up of the previous two ingredients: assume a real linear
combination vanishes, pair both sides with the conjugate of an arbitrary
holomorphic 1-form, apply the bilinear identity, and use Hermitian
positivity to force every coefficient to be zero. Why bother? It's the
rigorous statement that ``the period vectors fill out the full real
space.'' Equivalently, their integer combinations form a \emph{full
lattice} — and that lattice is what quotients to give the Jacobian its
compactness.
\end{layman}

\begin{proofsketch}
Suppose a real linear combination is zero. Pair both sides with
\(\bar\omega\) for an arbitrary basis \((\omega_i)\) of \(\HX\); the
bilinear identity together with Hermitian positivity forces all
coefficients to vanish.
\end{proofsketch}

\begin{theorem}[Classical input: Riemann bilinear relations (umbrella)]
\label{input:riemann-bilinear}
\uses{thm:bilinear-from-stokes,thm:hermitian-positivity,thm:period-vectors-full-real-rank}
\leanok
The classical Riemann bilinear relations and Hermitian positivity, with
their consequence that the \(2g\) symplectic-basis period vectors are
\(\R\)-linearly independent.
\lean{JacobianChallenge.Periods.period_vectors_linearIndependent_of_symplectic}
\end{theorem}

\begin{layman}
We package three deep facts as one classical input: the bilinear
identity expressing a wedge integral as a sum of period products,
Hermitian positivity of the self-pairing of a holomorphic 1-form, and
the consequence that the \(2g\) symplectic-basis period vectors are
linearly independent over the reals. Each is a tiny gem on its own,
but together they form the analytical core of the period-lattice
route. Why bother? The Jacobian's definition needs ``the integer span
of the period vectors is a full lattice in the dual space.'' That
single statement is formally a corollary of the umbrella, and naming
the umbrella keeps downstream constructions (Jacobian, Abel–Jacobi,
push-pull) from depending on the internal pieces.
\end{layman}

\begin{proof}[Proof of Theorem~\ref{thm:period-lattice}]\leanok
Choose a \(\C\)-basis \(\omega_1,\dots,\omega_g\) of \(\HX\), so
\(\HdualX\cong\C^g\). The images of the \(2g\) symplectic cycles under
\(\Pi_X\) are the columns of the period matrix.
By Theorem~\ref{thm:period-vectors-full-real-rank}, these \(2g\) vectors
are linearly independent over \(\R\). Since \(\HdualX\) has real
dimension \(2g\), their \(\Z\)-span is a full lattice.
\end{proof}

\begin{formalizationnote}
The current project has basis-aligned period subgroup files. The deepest
remaining content is the full-rank statement, routed through declarations such
as \code{periodSubgroup_isZLattice},
\code{periodSubgroup_spans_real}, and
\code{period_vectors_linearIndependent_of_symplectic}.
\end{formalizationnote}
